% بخش: روش‌ها
% فصل: اثبات‌ها
% بخش فرعی: الگوهای استنتاج

\documentclass[../../../include/open-logic-section]{subfiles}

\begin{document}

\olfileid{mth}{prf}{pat}

\olsection{الگوهای استنتاج}

اثبات‌ها از استنتاج‌های منفرد ساخته می‌شوند. هنگامی که استنتاجی
انجام می‌دهیم، معمولاً آن را با واژه‌هایی مانند «پس»، «بنابراین» یا
«ازاین‌رو» نشان می‌دهیم. استنتاج اغلب بر یک یا دو واقعیتی تکیه دارد
که پیشاپیش در اثبات در اختیار داریم؛ شاید چیزی باشد که فرض کرده‌ایم
یا چیزی که پیش‌تر با استنتاجی نتیجه گرفته‌ایم. برای روشنی، می‌توانیم
این مطالب را برچسب بزنیم و هنگام استنتاج مشخص کنیم از چه گزاره‌های
دیگری استفاده می‌کنیم. استنتاج اغلب توضیحی نیز دربردارد دربارهٔ
اینکه نتیجهٔ تازه \emph{چرا} از مطالب پیشین به دست می‌آید. چند
الگوی رایج استنتاج در اثبات‌ها بسیار به کار می‌روند؛ برخی را در
ادامه بررسی می‌کنیم. بعضی الگوهای استنتاج، مانند اثبات با استقرا،
پیچیده‌ترند (و بعداً بررسی خواهند شد).

پیش‌تر یک الگوی استنتاج را بررسی کرده‌ایم: بازکردن یا به‌کاربردن یک
تعریف. هنگامی که تعریفی را باز می‌کنیم، صرفاً مطلبی را که شامل
عبارت تعریف‌شونده است با استفاده از عبارت تعریف‌کننده بازمی‌گوییم. برای نمونه، فرض کنید
در جریان یک اثبات پیشاپیش $D = E$ را برقرار کرده‌ایم (الف). آن‌گاه
می‌توانیم تعریف $=$ را برای مجموعه‌ها به کار ببریم و نتیجه بگیریم:
«پس بنا بر تعریف و از (الف)، هر !!{element} از~$D$ !!a{element}
از~$E$ است و برعکس.»

کمی گیج‌کننده اینکه توجیه یک استنتاج را اغلب نه هنگام انجام آن،
بلکه پیش از آن می‌نویسیم. فرض کنید هنوز $D = E$ را اثبات نکرده‌ایم،
اما می‌خواهیم اثباتش کنیم. اگر $D = E$ نتیجهٔ مطلوب باشد، می‌توانیم
این هدف را نیز با به‌کاربردن تعریف بازگوییم: برای اثبات $D = E$
باید اثبات کنیم هر !!{element} از~$D$ !!a{element} از~$E$ است و
برعکس. پس اثبات ما چنین صورتی خواهد داشت: (الف) اثبات کنید هر
!!{element} از~$D$ !!a{element} از~$E$ است؛ (ب) هر !!{element}
از~$E$ !!a{element} از~$D$ است؛ (ج) پس از (الف) و (ب)، بنا بر
تعریف $=$، داریم $D = E$. اما معمولاً اثبات را این‌گونه نمی‌نویسیم.
در عوض، شاید چیزی مانند متن زیر بنویسیم:
\begin{quote}
می‌خواهیم $D = E$ را نشان دهیم. بنا بر تعریف~$=$، این کار به
نشان‌دادن این مطلب فروکاسته می‌شود که هر !!{element} از~$D$
!!a{element} از~$E$ است و برعکس.

(الف) \dots (اثباتی که هر !!{element} از~$D$ !!a{element}
از~$E$ است) \dots

(ب) \dots (اثباتی که هر !!{element} از~$E$ !!a{element}
از~$D$ است) \dots
\end{quote}

\subsection{استفاده از عطف}

شاید ساده‌ترین الگوی استنتاج، نتیجه‌گرفتن یکی از عطف‌شونده‌های یک
عطف باشد. به بیان دیگر، اگر فرض یا پیش‌تر اثبات کرده باشیم که $p$
و~$q$، حق داریم $p$ را نتیجه بگیریم (و نیز $q$ را). این استنتاج
چنان بنیادی است که اغلب ذکر نمی‌شود. برای نمونه، پس از آن‌که تعریف
$D = E$ را باز کردیم، برقرار کرده‌ایم که هر !!{element} از~$D$
!!a{element} از~$E$ است و برعکس. از این مطلب می‌توانیم نتیجه بگیریم
که هر !!{element} از~$E$ !!a{element} از~$D$ است (این همان بخش
«برعکس» است).

\subsection{اثبات یک عطف}

گاهی آنچه از شما خواسته می‌شود اثبات کنید صورت یک عطف را دارد؛ از
شما می‌خواهند «$p$ و $q$ را اثبات کنید». در این حالت فقط باید دو
کار انجام دهید: $p$ را اثبات کنید و سپس $q$ را. می‌توانید اثبات را
به دو بخش تقسیم کنید و برای روشنی به آن‌ها برچسب بزنید. هنگام
نوشتن یادداشت‌های نخست شاید بالای صفحه بنویسید «(۱) $p$ را اثبات
کنید» و میانهٔ صفحه بنویسید «(۲) $q$ را اثبات کنید». (البته شاید
به‌صراحت از شما نخواهند عطفی را اثبات کنید، اما دریابید که اثباتتان
به اثبات یک عطف نیاز دارد. برای نمونه، اگر از شما بخواهند
$D = E$ را اثبات کنید، پس از بازکردن تعریف~$=$ درمی‌یابید که باید
اثبات کنید: هر !!{element} از~$D$ !!a{element} از~$E$
\emph{و} هر !!{element} از~$E$ !!a{element} از~$D$ است.)

\subsection{اثبات یک فصل}

وقتی آنچه اثبات می‌کنید صورت یک فصل دارد (یعنی گزاره‌ای به صورت
«$p$ یا $q$» است)، کافی است نشان دهید یکی از فصل‌شونده‌ها صادق است.
اما تقریباً هرگز پیش نمی‌آید که یکی از آن دو فصل‌شونده بی‌درنگ از
فرض‌های قضیه نتیجه شود. اغلب، فرض‌های قضیه خود فصلی‌اند؛ یا نشان
می‌دهید همهٔ چیزهای نوعی معین یکی از دو ویژگی را دارند، اما بعضی
ویژگی نخست و بعضی دیگر ویژگی دوم را دارند. اینجاست که برهان به
حالت‌ها سودمند است (پایین‌تر را ببینید).

\subsection{اثبات شرطی}

بسیاری از قضایایی که با آن‌ها روبه‌رو می‌شوید صورت شرطی دارند (یعنی
باید نشان دهید اگر $p$ برقرار باشد، آن‌گاه $q$ نیز صادق است).
سامان‌دادن این موارد خوب و آسان است: کافی است مقدم شرطی (در اینجا
$p$) را فرض کنید و نتیجهٔ~$q$ را از آن اثبات کنید. پس اگر قضیهٔ شما
می‌گوید «اگر $p$، آن‌گاه $q$»، اثبات را با «$p$ را فرض کنید» آغاز
می‌کنید و در پایان باید~$q$ را اثبات کرده باشید.

شرطی‌ها ممکن است به شیوه‌های متفاوتی بیان شوند. پس به‌جای «اگر
$p$، آن‌گاه $q$»، ممکن است قضیه‌ای بگوید «$p$ تنها اگر $q$»،
«$q$ اگر $p$» یا «$q$، مشروط بر $p$». همهٔ این‌ها معنای یکسانی
دارند و مستلزم فرض‌کردن $p$ و اثبات~$q$ از آن فرض‌اند. به یاد آورید
که یک دوشرطی («$p$ اگر و تنها اگر $q$») در واقع دو شرطی است که کنار
هم قرار گرفته‌اند: اگر $p$، آن‌گاه $q$؛ و اگر $q$، آن‌گاه~$p$.
پس فقط باید دو نمونه اثبات شرطی انجام دهید، یکی برای شرطی نخست و
دیگری برای شرطی دوم. بااین‌حال، گاهی می‌توان گزاره‌ای «اگر و تنها
اگر» را با زنجیرکردن چند گزارهٔ «اگر و تنها اگر» دیگر اثبات کرد،
به‌طوری‌که با «$p$» آغاز و با «$q$» پایان دهید؛ اما در این صورت
باید مطمئن شوید هر گام واقعاً یک «اگر و تنها اگر» است.

\subsection{ادعاهای کلی}

استفاده از یک ادعای کلی ساده است: اگر چیزی دربارهٔ هر شیء صادق باشد،
دربارهٔ هر شیء خاص نیز صادق است. برای نمونه، اگر فرض اثبات شما
$A \subseteq B$ باشد، این (با بازکردن تعریف~$\subseteq$) بدان معناست
که برای هر $x \in A$، داریم $x \in B$. بنابراین، اگر از پیش بدانید
$z \in A$، می‌توانید $z \in B$ را نتیجه بگیرید.

اثبات یک ادعای کلی شاید کمی دشوار به نظر برسد. معمولاً این گزاره‌ها
صورت زیر را دارند: «اگر $x$ ویژگی~$P$ را دارد، آنگاه ویژگی~$Q$ را
نیز دارد» یا «همهٔ $P$ها $Q$اند». البته ممکن است گزاره دقیقاً در
این قالب نگنجد و تشخیص دقیقِ آنچه باید اثبات کنید کمی تمرین بخواهد.
اما اغلب باید نشان دهیم همهٔ اشیای دارای یک ویژگی، ویژگی معین دیگری
نیز دارند.

راه اثبات یک ادعای کلی این است که برای چیزهای دارای ویژگی نخست نام
یا متغیری معرفی کنیم و سپس نشان دهیم آن‌ها ویژگی دوم را نیز دارند.
می‌توانیم چنین بگوییم: برای اثبات چیزی دربارهٔ \emph{همهٔ} $P$ها،
باید آن را دربارهٔ یک $P$ \emph{دلخواه} اثبات کنید. نام معرفی‌شده،
نام یک $P$ دلخواه است. معمولاً از حروف منفرد برای نام‌گذاری این
اشیای دلخواه استفاده می‌کنیم و حروف از قراردادهایی پیروی می‌کنند؛
برای نمونه، $n$ را برای اعداد طبیعی، $!A$ را برای !!{formula}ها،
$A$ را برای مجموعه‌ها، $f$ را برای توابع، و جز این‌ها به کار
می‌بریم.

نکتهٔ ظریف، حفظ کلیت در سراسر اثبات است. با فرض اینکه شیئی دلخواه
(«$x$») ویژگی~$P$ را دارد آغاز می‌کنید و (فقط بر پایهٔ تعریف‌ها یا
آنچه اجازه دارید فرض کنید) نشان می‌دهید $x$ ویژگی~$Q$ را دارد. چون
مشخص نکرده‌اید $x$ دقیقاً چیست، جز اینکه ویژگی $P$ را دارد، می‌توانید
اعلام کنید هر چیزی با ویژگی $P$، ویژگی~$Q$ را دارد. خلاصه، $x$
جانشینِ \emph{همهٔ} چیزهای دارای ویژگی~$P$ است.

\begin{prop}
  برای همهٔ مجموعه‌های $A$ و $B$،
  $A \subseteq A \cup B$.
\end{prop}

\begin{proof}
  بگذارید $A$ و $B$ مجموعه‌هایی دلخواه باشند. می‌خواهیم نشان دهیم
  $A \subseteq A \cup B$. بنا بر تعریف $\subseteq$، این بدان
  فروکاسته می‌شود که برای هر $x$، اگر $x \in A$، آن‌گاه
  $x \in A \cup B$. پس بگذارید $x \in A$ یک !!{element} دلخواه
  از~$A$ باشد. باید نشان دهیم $x \in A \cup B$. چون $x \in A$،
  داریم $x \in A$ یا $x \in B$. بنابراین
  $x \in \Setabs{x}{x \in A \lor x \in B}$. اما این، بنا بر تعریف
  $\cup$، یعنی $x \in A \cup B$.
\end{proof}


\subsection{برهان به حالت‌ها}

فرض کنید فصلی را به‌عنوان فرض یا نتیجه‌ای ازپیش‌برقرارشده دارید؛
فرض یا اثبات کرده‌اید که $p$ یا $q$ صادق است. می‌خواهید $r$ را
اثبات کنید. این کار را در دو گام انجام می‌دهید: نخست صدق $p$ را فرض
و~$r$ را اثبات می‌کنید؛ سپس صدق $q$ را فرض و بار دیگر~$r$ را اثبات
می‌کنید. این روش کار می‌کند، زیرا فرض کرده‌ایم یا می‌دانیم یکی از
دو بدیل برقرار است. دو گام نشان می‌دهند که هریک از آن دو برای صدق
~$r$ کافی است. (اگر هر دو صادق باشند، نه یک، بلکه دو دلیل برای صدق
$r$ داریم. لازم نیست جداگانه اثبات کنیم که با فرض هم‌زمان $p$ و
~$q$، گزارهٔ $r$ صادق است.) برای نشان‌دادن کاری که انجام می‌دهیم،
اعلام می‌کنیم «حالت‌ها را جدا می‌کنیم». برای نمونه، فرض کنید
می‌دانیم $x \in B \cup C$. مجموعهٔ $B \cup C$ به‌صورت
$\Setabs{x}{x \in B \text{ یا } x \in C}$ تعریف می‌شود. به بیان
دیگر، بنا بر تعریف، $x \in B$ یا $x \in C$. برای اثبات
$x \in A$ از این مطلب، نخست $x \in B$ را فرض و از آن
$x \in A$ را اثبات می‌کنیم؛ سپس $x \in C$ را فرض و دوباره
$x \in A$ را اثبات می‌کنیم. زیر فرض می‌نویسید «حالت‌ها را جدا
می‌کنیم»، سپس «حالت (۱): $x \in B$» را پایین‌تر و «حالت (۲):
$x \in C$» را در نیمهٔ صفحه می‌نویسید. آنگاه نیمهٔ بالا و نیمهٔ
پایین صفحه را کامل می‌کنید.

برهان به حالت‌ها به‌ویژه هنگامی سودمند است که خودِ آنچه اثبات
می‌کنید فصلی باشد. این مثال ساده را ببینید:

\begin{prop}
فرض کنید $B \subseteq D$ و $C \subseteq E$. آن‌گاه
$B \cup C \subseteq D \cup E$.
\end{prop}

\begin{proof}
  فرض کنید (الف) $B \subseteq D$ و (ب) $C \subseteq E$. بنا بر
  تعریف، هر $x \in B$ در $D$ نیز هست، یعنی $x \in D$ (ج)؛ و هر
  $x \in C$ در $E$ نیز هست، یعنی $x \in E$ (د). برای نشان‌دادن
  $B \cup C \subseteq D \cup E$ باید نشان دهیم اگر
  $x \in B \cup C$، آن‌گاه $x \in D \cup E$ (بنا بر تعریف
  $\subseteq$). بنا بر تعریف~$\cup$،
  $x \in B \cup C$ اگر و تنها اگر $x \in B$ یا $x \in C$. به همین
  ترتیب، $x \in D \cup E$ اگر و تنها اگر $x \in D$ یا $x \in E$.
  پس باید نشان دهیم: برای هر $x$، اگر $x \in B$ یا $x \in C$،
  آن‌گاه $x \in D$ یا $x \in E$.

  \begin{quote}
  تا اینجا فقط تعریف‌ها را باز کرده‌ایم!{} حکم خود را بدون
  $\subseteq$ و $\cup$ بازصورت‌بندی کرده‌ایم و اکنون باید ادعایی
  کلی و شرطی را اثبات کنیم. بنا بر بحث بالا، این کار با فرض‌کردن
  اینکه $x$ چیزی است که بخش «اگر» درباره‌اش صادق است انجام می‌شود؛
  سپس نشان می‌دهیم بخش «آنگاه» نیز صادق است. به بیان دیگر، فرض
  می‌کنیم $x \in B$ یا $x \in C$ و نشان می‌دهیم
  $x \in D$ یا $x \in E$.\footnote{این بند فقط کاری را که انجام
  می‌دهیم توضیح می‌دهد؛ بخشی از اثبات نیست و هنگام نوشتن اثبات‌های
  خود لازم نیست این‌همه وارد جزئیات شوید.}
  \end{quote}

  فرض کنید $x \in B$ یا $x \in C$. باید نشان دهیم $x \in D$ یا
  $x \in E$. حالت‌ها را جدا می‌کنیم.

  حالت ۱: $x \in B$. بنا بر (ج)، $x \in D$. پس $x \in D$ یا
  $x \in E$. (اینجا استنتاجِ بررسی‌شده در زیربخش پیشین را انجام
  داده‌ایم!)

  حالت ۲: $x \in C$. بنا بر (د)، $x \in E$. پس $x \in D$ یا
  $x \in E$.
 \end{proof}


\subsection{اثبات یک ادعای وجودی}

هنگامی که از شما می‌خواهند ادعایی وجودی را اثبات کنید، پرسش معمولاً
این صورت را دارد: «اثبات کنید $x$ای هست به‌طوری‌که
$\dots x \dots$»؛ یعنی اثبات کنید شیئی وجود دارد که ویژگیِ توصیف‌شده
با «$\dots x \dots$» را دارد. در این حالت باید شیئی مناسب مشخص کنید
و نشان دهید ویژگی لازم را دارد. این کار سرراست به نظر می‌رسد، اما
اثباتی از این نوع ممکن است دشوار باشد. معمولاً مستلزم
\emph{ساختن} یا \emph{تعریف‌کردن} یک شیء و اثبات این است که شیء
تعریف‌شده ویژگی لازم را دارد. یافتن شیء درست ممکن است دشوار باشد؛
اثبات ویژگی لازم شاید دشوار باشد؛ و گاهی حتی نشان‌دادن اینکه اصلاً
در تعریف یک شیء موفق شده‌اید نیز دشوار است!{}

معمولاً این را با مشخص‌کردن شیء می‌نویسید؛ برای نمونه، «بگذارید
$x$ برابر با \dots باشد» (که در آن \dots{} شیء مورد نظر را مشخص
می‌کند). شاید لازم باشد اثبات کنید \dots واقعاً شیئی موجود را توصیف
می‌کند؛ سپس نشان می‌دهید $x$ ویژگی~$Q$ را دارد. این مثال ساده را
ببینید.

\begin{prop}
  فرض کنید $x \in B$. آن‌گاه $A$ای وجود دارد به‌طوری‌که
  $A \subseteq B$ و $A \neq \emptyset$.
\end{prop}

\begin{proof}
  فرض کنید $x \in B$. بگذارید $A = \{x\}$.
  \begin{quote}
    اینجا مجموعهٔ~$A$ را با برشمردن !!{element}هایش تعریف کرده‌ایم.
    چون فرض می‌کنیم $x$ یک شیء است و همیشه می‌توانیم مجموعه‌ای را
    با برشمردن !!{element}هایش تشکیل دهیم، لازم نیست اینجا نشان دهیم
    که در تعریف مجموعهٔ~$A$ موفق شده‌ایم. بااین‌حال، هنوز باید نشان
    دهیم $A$ ویژگی‌های خواسته‌شده در حکم را دارد. بدون آن، اثبات
    کامل نیست!{}
  \end{quote}
  چون $x \in A$، داریم $A \neq \emptyset$.
  \begin{quote}
    این مطلب بر تعریف $A$ به‌صورت $\{x\}$ و دو واقعیت بدیهی تکیه
    دارد: $x \in \{x\}$ و $x \notin \emptyset$.
  \end{quote}
  چون $x$ تنها !!{element} از~$\{x\}$ است و $x \in B$، هر
  !!{element} از~$A$ !!a{element} از~$B$ نیز هست. بنا بر تعریف
  $\subseteq$، داریم $A \subseteq B$.
\end{proof}

\subsection{استفاده از ادعاهای وجودی}

فرض کنید می‌دانید ادعایی وجودی صادق است (آن را اثبات کرده‌اید یا
فرضی است که می‌توانید به کار برید)؛ مثلاً «برای یک~$x$،
$x \in A$» یا «$x \in A$ای وجود دارد». اگر بخواهید آن را در
اثبات به کار ببرید، می‌توانید وانمود کنید برای یکی از چیزهایی که
فرضتان وجودشان را اعلام می‌کند نامی دارید. چون $A$ دست‌کم یک چیز
دارد، چیزهایی هستند که آن نام ممکن است به آن‌ها اشاره کند. البته
شاید نتوانید یکی را مشخص کنید یا بیش از اینکه در~$A$ است توصیفش
کنید. اما برای مقصود اثبات می‌توانید وانمود کنید آن را مشخص کرده‌اید
و نامی بدان بدهید. مهم است نامی برگزینید که پیش‌تر به کار نبرده‌اید
(و در فرض‌هایتان نیامده است)، وگرنه ممکن است اشکالی پیش آید. در
اثبات این کار را با گذر از «برای یک $x$، $x \in A$» به «بگذارید
$a \in A$» نشان می‌دهید. اکنون می‌توانید دربارهٔ~$a$ استدلال کنید،
فرض‌های دیگر را به کار ببرید، و جز این‌ها، تا به نتیجه‌ای مانند
$p$ برسید. اگر $p$ دیگر از~$a$ نام نبرد، $p$ از فرض
$a \in A$ مستقل است و نشان داده‌اید که فقط از فرضِ
«برای یک $x$، $x \in A$» نتیجه می‌شود.

\begin{prop}
اگر $A \neq \emptyset$، آن‌گاه $A \cup B \neq \emptyset$.
\end{prop}

\begin{proof}
  فرض کنید $A \neq \emptyset$. پس برای یک $x$، داریم $x \in A$.
  \begin{quote}
    اینجا نخست فقط فرض حکم را بازگفتیم. این فرض، یعنی
    $A \neq \emptyset$، ادعایی وجودی را پنهان می‌کند که تنها با
    بازکردن چند تعریف بدان می‌رسیم. تعریف $=$ به ما می‌گوید
    $A = \emptyset$ اگر و تنها اگر هر $x \in A$ در
    $\emptyset$ نیز باشد، یعنی $x \in \emptyset$، و هر
    $x \in \emptyset$ در $A$ نیز باشد، یعنی $x \in A$. با نقیض‌کردن
    دو سو به دست می‌آوریم:
    $A \neq \emptyset$ اگر و تنها اگر یا یک $x \in A$ در
    $\emptyset$ نباشد، یا یک $x \in \emptyset$ در $A$ نباشد. چون
    هیچ چیز در $\emptyset$ نیست، فصل‌شوندهٔ دوم هرگز نمی‌تواند صادق
    باشد؛ و «$x \in A$ و $x \notin \emptyset$» فقط به
    $x \in A$ فروکاسته می‌شود. پس
    $A \neq \emptyset$ اگر و تنها اگر برای یک $x$،
    $x \in A$. این یک ادعای وجودی است. اکنون با معرفی نامی برای یکی
    از !!{element}های~$A$، آن ادعای وجودی را به کار می‌بریم:
  \end{quote}
  بگذارید $a \in A$.
  \begin{quote}
    اکنون نامی برای یکی از چیزهای در~$A$ معرفی کرده‌ایم. استدلال
    دربارهٔ~$a$ را ادامه می‌دهیم، اما مراقبیم فقط $a \in A$ را فرض
    کنیم و نه هیچ چیز دیگر:
  \end{quote}
  چون $a \in A$، بنا بر تعریف~$\cup$ داریم
  $a \in A \cup B$. پس برای یک $x$، $x \in A \cup B$؛ یعنی
  $A \cup B \neq \emptyset$.
  \begin{quote}
    در گام آخر از «$a \in A \cup B$» به «برای یک $x$،
    $x \in A \cup B$» گذشتیم. عبارت دوم دیگر از $a$ نام نمی‌برد،
    پس می‌دانیم «برای یک $x$، $x \in A \cup B$» تنها از
    «برای یک $x$، $x \in A$» نتیجه می‌شود. اما این یعنی
    $A \cup B \neq \emptyset$.
  \end{quote}
\end{proof}

شاید تمرین خوبی باشد که متغیرهای مقید مانند «$x$» را، همان‌گونه که
ما انجام دادیم، از نام‌های فرضی مانند $a$ جدا نگه دارید. اما در عمل
اغلب چنین نمی‌کنیم و فقط $x$ را به کار می‌بریم، مثلاً:
\begin{quote}
فرض کنید $A \neq \emptyset$؛ یعنی $x \in A$ای وجود دارد. بنا بر
تعریف $\cup$، داریم $x \in A \cup B$. پس
$A \cup B \neq \emptyset$.
\end{quote}
اما هنگام انجام این کار باید بسیار مراقب باشید که برای ادعاهای
وجودی متفاوت از $x$ها و $y$های متفاوت استفاده کنید. برای نمونه،
متن زیر اثبات درستی از «اگر $A \neq \emptyset$ و
$B \neq \emptyset$، آن‌گاه $A \cap B \neq \emptyset$» نیست (و خود
این گزاره نیز صادق نیست).
\begin{quote}
فرض کنید $A \neq \emptyset$ و $B \neq \emptyset$. پس برای یک $x$،
$x \in A$؛ و نیز برای یک $x$، $x \in B$. چون $x \in A$ و
$x \in B$، بنا بر تعریف~$\cap$ داریم $x \in A \cap B$. پس
$A \cap B \neq \emptyset$.
\end{quote}
آیا می‌توانید گام نادرست را بیابید و توضیح دهید چرا نتیجه برقرار
نیست؟

\end{document}
